%%%%%%%%%%%%%%%%%%%%%%%%%%%%%%%%%%%%%%%%%
% Medium Length Professional CV
% LaTeX Template
% Version 3.0 (December 17, 2022)
%
% This template originates from:
% https://www.LaTeXTemplates.com
%
% Author:
% Vel (vel@latextemplates.com)
%
% Original author:
% Trey Hunner (http://www.treyhunner.com/)
%
% License:
% CC BY-NC-SA 4.0 (https://creativecommons.org/licenses/by-nc-sa/4.0/)
%
%%%%%%%%%%%%%%%%%%%%%%%%%%%%%%%%%%%%%%%%%

%----------------------------------------------------------------------------------------
%  PACKAGES AND OTHER DOCUMENT CONFIGURATIONS
%----------------------------------------------------------------------------------------

\documentclass[
  %a4paper, % Uncomment for A4 paper size (default is US letter)
  11pt, % Default font size, can use 10pt, 11pt or 12pt
]{resume} % Use the resume class

\usepackage{ebgaramond} % Use the EB Garamond font
\usepackage{ulem}
\usepackage{hyperref}

\AtBeginDocument{
  \let\oldhref=\href
  \renewcommand\href[2]{\dotuline{\oldhref{#1}{#2}}}
}

\pagestyle{plain}
%------------------------------------------------

\name{Arman Babaei} % Your name to appear at the top

% You can use the \address command up to 3 times for 3 different
% addresses or pieces of contact information
% Any new lines (\\) you use in the \address commands will be
% converted to symbols, so each address will appear as a single line.

% \address{123 Broadway \\ City, State 12345} % Main address

% \address{123 Pleasant Lane \\ City, State 12345} % A secondary
% address (optional)

\address{+98~$\cdot$~919~$\cdot$~043~$\cdot$~7966
\\ ar.babaei@sharif.edu} % Contact information

%----------------------------------------------------------------------------------------

\begin{document}

%\vspace{20pt}

%----------------------------------------------------------------------------------------
%  EDUCATION SECTION
%----------------------------------------------------------------------------------------

\begin{rSection}{Education}
  \vspace{5pt}
  \begin{rSubsection}{Sharif University of Technology}{September 2024
    -- Expected July 2026}{M.S. in Computer Engineering (Networks)}{}{}
  \item Current GPA: 18.60 (/20)
  \item TOEFL Score: 113/120
  \end{rSubsection}

  %------------------------------------------------

  \begin{rSubsection}{Sharif University of Technology}{September 2019
    -- July 2024}{B.S. in Computer Engineering}{}{}
  \item Overall GPA: 19.48 (/20)
  \item Rank: 15/190
  \end{rSubsection}

  %------------------------------------------------

  \begin{rSubsection}{Allameh Helli Highschool}{2016 -- 2019}{}{}
  \item Physics and Mathematics
  \item Overall GPA: 19.74 (/20)
  \end{rSubsection}
\end{rSection}

%----------------------------------------------------------------------------------------
%  HONORS
%----------------------------------------------------------------------------------------
\begin{rSection}{Honors}
  \vspace{3pt}
  \begin{tabular}{@{} >{\bfseries}l @{\hspace{6ex}} l @{}}
    \textbf{2024} & Ranked \textbf{\#4} in Iran's Master University
    Entrance Exam. (System and network subgroup) \\
    \textbf{2019} & Ranked top \textbf{0.01\%} in Iran's National
    University Entrance Exam. \\
    \textbf{2018} & \textbf{Bronze} medalist of Iran National
    Olympiad in Informatics. (\textbf{INOI}) \\
    % \textbf{2014} & \textbf{3rd} place in Sharif Mobile Programming Marathon\\
  \end{tabular}
\end{rSection}

% \vspace{15pt}
%----------------------------------------------------------------------------------------
%  RESEARCH EXPERIENCE SECTION
%----------------------------------------------------------------------------------------

\begin{rSection}{Research Experience}
  \vspace{5pt}

  \begin{rSubsection}{Sharif University of Technology}{August 2025 -
    Present}{L4S/AQM Prototyping for Video Streaming}{Tehran, Iran}
  \item Prototyping L4S-style low-latency AQM using DualPI2 in P4\_16
    (BMv2 \texttt{v1model}) with per-egress-port register state;
    studying how ingress traffic classification interacts with egress
    AQM behavior; and comparing against Linux DualPI2
    implementations. Exploring different queuing mechanisms to
    improve latency and QoS for volumetric video streaming.
  \end{rSubsection}

  \begin{rSubsection}{Stony Brook University}{October 2024 -
    Present}{Adaptive Leader Election for PBFT}{remote}
  \item Designing an adaptive leader-election module for a PBFT-style
    protocol with a component-based architecture (PBFT core,
    ViewChanger, LeaderElection) and clear interfaces for state
    reporting and control. Developing a learning-based loop that
    aggregates node-reported execution progress via Byzantine-robust
    consensus and median estimation to adapt leader-election decisions.
  \end{rSubsection}

  \begin{rSubsection}{EPFL (École Polytechnique Fédérale de
    Lausanne)}{July 2023 - September 2023}{Fairness of Incentive
    Mechanism in Distributed Storage Systems}{Lausanne, Switzerland}
  \item Working in \textbf{Bryan Ford}'s lab
    (\href{https://www.epfl.ch/labs/dedis/}{DEDIS}) supervised by
    Vero Estrada-Galiñanes on income distribution and fairness in
    the Swarm storage network. I worked on a mathematical model that
    enables analyzing network properties
    without simulation, and I added the analysis for different node
    behaviors and nodes' churn to a simulator written
    in Go. I also developed and evaluated a network shuffling method
    which improved the fairness. (preprint
    available at \href{https://arxiv.org/abs/2307.02231}{arXiv})
  \end{rSubsection}

  %------------------------------------------------

  \pagebreak

  \begin{rSubsection}{ISTA (Institute of Science and Technology
    Austria)}{July 2022 - February 2023}{Parallel Execution in
    Blockchain Protocols}{Vienna, Austria}
  \item Working in the
    \href{https://ist.ac.at/en/research/kokoris-group/}{SPiDerS lab}
    of the Institute of Science and Technology Austria under
    supervision of \textbf{Lefteris Kokoris-Kogias}, responsible for
    reproducing the
    previous results and setting up a test bench and other
    implementation-needed issues, such as analyzing the
    current workload of commercial blockchains. Contributing to the
    discovery of potential aids that can be used to
    improve performance. (preprint
    available at \href{https://arxiv.org/abs/2401.14278}{arXiv})
  \end{rSubsection}

  %------------------------------------------------

  %   \begin{rSubsection}{Sharif University of Technology}{December 2021
  %     - February 2022}{Resource Allocation in Fog Computing}{Tehran, Iran}
  %   \item Researching under supervision of Ms Hajar (Maryam) Siar, PhD
  %     student of \textbf{Dr Mohammad Izadi}, about fog computing
  %     challenges in resource allocation and available techniques.
  %     This Consists of studying an overview of challenges in cloud,
  %     edge, and fog computing, and also approaches used
  %     for tackling the resource allocation problem and trying to come
  %     up with ideas in this area.
  %   \end{rSubsection}

  %------------------------------------------------

  %   \vspace{10pt}

  \begin{rSubsection}{Sharif University of Technology}{May 2021 -
    July 2021}{Distributed Shortest Path in CONGEST Model}{Tehran, Iran}
  \item Analyzing upper and lower bounds for finding
    single/multi-source shortest path in CONGEST model and
    trying to come up with new ideas for lowering the upper bound of
    single source shortest path problem.
  \end{rSubsection}

\end{rSection}

% \pagebreak

% \vspace{10pt}

%----------------------------------------------------------------------------------------
%  PUBLICATIONS SECTION
%----------------------------------------------------------------------------------------

\begin{rSection}{Publications}
  \vspace{5pt}

  \textbf{Pythia: Supercharging Parallel Smart Contract Execution to Guide Stragglers and Full Nodes to Safety} \hfill
  {\href{https://arxiv.org/abs/2401.14278}{SIB} -- }\\
  \phantom{} \hfill 2024
  \smallskip

  \textbf{Altruism, reciprocity, and tokens to reward forwarding data: Is that fair?} \hfill
  {\href{https://arxiv.org/abs/2307.02231}{ICBC} -- 2024}
  \smallskip
\end{rSection}

%----------------------------------------------------------------------------------------
%  WORK EXPERIENCE SECTION
%----------------------------------------------------------------------------------------
% \vspace*{15pt}
\begin{rSection}{Work Experience}
  \vspace{5pt}
  \begin{rSubsection}{Tapsi}{September 2024 -- Present}{Proficient Platform Engineer}{Tehran, Iran}
      \item \textbf{Apr'26 -- Present} Conducting platform R\&D for reliability, deployment safety, observability, and dev experience.
      \item \textbf{Sep'25 -- Apr'26} Improved backend reliability through monitoring, incident analysis, and operations.
      \item \textbf{Sep'24 -- Sep'25} Built production passenger--driver matching features for a ride-hailing system.
  \end{rSubsection}

  %------------------------------------------------

  \begin{rSubsection}{Undisclosed Hedge Fund}{September 2023 -
    February 2024}{Full-stack Developer}{remote}
  \item Developing a job scheduler web service for intra-company tasks.
  \end{rSubsection}

  %------------------------------------------------

  \begin{rSubsection}{Mahsan co.}{September 2020 - July
    2022}{C++/Python Developer}{Tehran, Iran}
  \item Python Developer, Developing a domain-specific language for
    end-to-end testing.
  \item C++ Developer, Developing a parallel network relaying software.
  \item Scrum Master, leading a team of five.
  \end{rSubsection}

  %------------------------------------------------

  \begin{rSubsection}{Allameh Helli Highschool}{June 2019 - October
    2020}{Algorithm Design Teacher}{Tehran, Iran}
  \item Teaching algorithm design to students with the aim of
    competing in the Iran National Olympiad in Informatics.
  \end{rSubsection}

\end{rSection}

% \vspace{10pt}

%----------------------------------------------------------------------------------------
%  TEACHING ASSISTANT
%----------------------------------------------------------------------------------------

\begin{rSection}{Teaching Assistant}
  \vspace{5pt}
  \begin{tabular}{@{} >{\bfseries}l @{\hspace{6ex}} l @{\hspace{6ex}} l @{}}
    Operating Systems & F'23 & Responsible for designing the
    assignments of the course. \\
    Game Theory & S'23 & Designed and evaluated three practical assignments.\\
    Advanced Programming & S'22,S'23 & Responsible for designing and
    coordinating the course project. \\
    Fundamentals of the Programming & F'22 & Responsible for
    designing and coordinating the course project. \\
    Artificial Intelligence & S'22,F'22 & Helped in designing and
    grading assignments. \\
    Data Structures and Algorithms & S'22 & Responsible for designing
    midterm and final exams. \\
    Logic Design & S'22 & Designed and examined one of the
    assignments of the semester.
  \end{tabular}
\end{rSection}

% \vspace{10pt}

% \begin{rSection}{Influential Courses}
%   \vspace{5pt}
%   \centering
%   \begin{tabular}{@{} l @{\hspace{4ex}} l @{\hspace{4ex}} l @{}}
%     Distributed Systems & Theory of Machines and Languages & Design
%     of Algorithms \\
%     Game Theory & Machine Learning & Multicore processing \\
%   \end{tabular}\par
% \end{rSection}

%----------------------------------------------------------------------------------------
%  EXAMPLE SECTION
%----------------------------------------------------------------------------------------

%\begin{rSection}{Section Name}

%Section content\ldots

%\end{rSection}

%----------------------------------------------------------------------------------------

\end{document}
